\section{Continuity with version 1 and the role of the literature}
\label{app:v1-record}

This revision contains the mathematical arguments of the original
preprint \cite{v1} within the expanded exposition. The square
identity is Lemma~\ref{lem:square}; its exact finite-dimensional
duality statement uses the corrected extremizer
\(\beta_h=\overline{C_h}/\|C\|_2\).
Proposition~\ref{prop:v1-terminal} retains the original
terminal-\(\eta\) criterion. The maximal argument appears in
Theorem~\ref{thm:bessel} on a common support, with its
\((K+1)^2\) cost displayed before absorbing logarithms into
an arbitrarily small power. The padding covers the entire
support: for coefficients on \((N,2N]\), use
\(2N\le2^K<4N\). No assertion about arbitrary translated
intervals follows from a spectral bound tied to the size of
the indices.

Proposition~\ref{prop:uniform} strengthens the original
coefficient-uniform failure theorem to an exact maximum.
Corollary~\ref{cor:v1-power-gap} and Figure~\ref{fig:v1-gap}
retain its particular \(9/4\), \(15/8\) comparison.
Lemma~\ref{lem:v1-square-divisor} retains the square-divisor
decomposition and displays every correlation stratum. The
original conditional RH conclusion is included, and strengthened
to a conditional Dirichlet-GRH implication, in
Theorem~\ref{thm:grh}. The elementary RH criterion and its
reverse direction are given in Theorem~\ref{thm:equiv} and
Appendix~\ref{sec:perron}. These are implications, not proofs
of their correlation hypotheses.

\subsection{The analytic context retained from the original paper}

The elementary identities sit within the standard analytic
number theory of \cite{apostol,iwaniec-kowalski,titchmarsh}.
The spectral input in Section~\ref{sec:sieve} belongs to the
large-sieve framework of
\cite{di,iwaniec-spectral,friedlander-iwaniec,bombieri,montgomery}.
The present maximal argument is an elementary consequence of a
fixed-support Bessel inequality; it does not supply a new spectral
input. Work on exceptional eigenvalues \cite{pascadi} addresses
a different refinement.

The role of signs is also central in M\"obius disjointness
\cite{sarnak-letter,bourgain-sarnak-ziegler} and in the correlation
theorems of \cite{mr-short,mrt,matomaki-radziwill-tao-signs,tao,
tao-teravainen,klurman,teravainen-binary}. The related settings
of Bohr sets, partial sums and Chowla variants are represented
by \cite{teravainen-walker,chavez,krishnamoorthy-chowla}.
These references are retained for context. None is taken here
as a proof of the maximal collar or of the square-root
prime-tail estimate. Classical zero results \cite{selberg},
finite-height verification \cite{pt}, and the disproof of the
Mertens conjecture \cite{odlyzko-teriele} likewise have their
own scopes. The character-sum viewpoint is represented by
\cite{granville-soundararajan}; the present character criteria
do not import an unstated uniform bound from that theory.

The recent Hexagon paper \cite{caceres} studies weighted-prime
margin criteria, reports finite verification, and analyzes
the obstruction to inferring a global bound from averaged
information. Its discussion of moving scales and exceptional
sets is pertinent to our kernel comparison. We use it as
related work, not as an input to any theorem here. In
particular, finite verification does not establish the bound
at every \(X\), and a criterion equivalent to RH does not
establish its required inequality.

\subsection{Funding and availability}
This research received no external funding. No empirical dataset
is used. The full research package contains the supplied v1 TeX
source and its original supplemental scripts under
\texttt{source-history/v1/}, together with their checksums and
a section-by-section content map. The original exponent check
and optional plotting script are preserved without alteration.
The latter requires NumPy and Matplotlib; it is not needed to
compile the manuscript, whose comparison figure is native
TikZ/PGFPlots. The new exact checks use Python's standard library.

\subsection{Ethics and interests}
The work is theoretical mathematics and involves no human
participants, animal subjects or clinical data requiring ethical
approval. The authors declare no conflicts of interest. The
computational and AI-assistance statement, including the
limited scope of the Lean verification, follows below.
